% ===== §1 Prelude: What Is Wealth? =====
\section{Prelude: What Is Wealth?}
\label{sec:prelude}

\noindent
Consider three scenes.

In the first, a family negotiates a bride-price. Gold changes hands---a quantity agreed upon after careful deliberation, reflecting the woman's family's assessment of what they are giving and the man's family's assessment of what they are receiving. The transaction is witnessed, documented, culturally sanctioned. When it is complete, both parties know exactly what has been transferred. The wealth is real, visible, verifiable.

In the second scene, two people sit across from each other after fifteen years of shared life. Between them, invisible to any accountant, lies something that neither owns individually and neither could enumerate: the accumulated texture of a life built together---the arguments and reconciliations, the ordinary Tuesday evenings and the extraordinary moments, the flowers noticed on paths and shared without calculation, the griefs endured together and the joys that arose from their \emph{between}. No document records this. No court could adjudicate its distribution. And yet both people, in some inarticulate way, know that it is the most significant thing either of them possesses.

In the third scene, that same couple, after years of exhaustion and the slow accumulation of unshared contingency, finds themselves strangers in the same house. The gold from the first scene---if there was any---remains. The bank accounts are unchanged. But something else has been depleted that cannot be measured in any currency, and its absence is felt as a poverty more complete than any material lack.

\emc{What is wealth?}

This question, asked in the context of intimate relations, turns out not to be an economic question at all. It is an ontological question---a question about what kinds of things exist, how they are generated, and who can possess them. The form of wealth that organises intimate relations determines whether the relational field can genuinely flourish, whether both parties can participate as genuine co-subjects of a shared life, or whether one or both are structurally reduced to exchangeable units within a logic of accumulation that is fundamentally alien to what intimate relations are for.

The paper that follows is an attempt to think carefully about this question. It proceeds from a conviction that the major traditions of wealth theory---the mercantilism that treats wealth as accumulation and exchange, the Marxist tradition that analyses wealth as the product of exploited labour, the capabilities approach that reconceives wealth as freedom, the gift economy that understands wealth as flow and social bond---all illuminate something real about wealth in intimate relations, and all miss something equally real. What they miss, this paper argues, is a form of wealth that is irreducible to any of their categories: a form of wealth that is generated through co-evolutionary activity in the relational field, that belongs to the \emph{between} rather than to either individual, and whose most fundamental property is that use does not deplete but may augment.

We call this \emb{Generative Relational Wealth} (GRW).

\begin{claim}[The central claim]
The wealth question in intimate relations is not an economic question but an ontological one: the form of wealth that organises intimate relations determines whether relational subjects can genuinely exist as co-evolving beings. GRW---wealth generated through co-evolutionary activity in the relational field---is the primary form of wealth in intimate relations. It is inexhaustible in the taking and inexhaustible in the using, because it is stored not in any material or symbolic medium but in the permanent structural modifications of the relational attractor landscape. Without understanding GRW, we cannot understand what it means for an intimate relation to flourish---or to fail.
\end{claim}

\emc{A word on method.} This paper is, by necessity, interdisciplinary to a degree that may initially seem excessive. The question of wealth in intimate relations requires political economy (to map the existing traditions and their assumptions), philosophy (to interrogate the ontological foundations), psychoanalysis (to understand the symbolic structures through which wealth and power operate in intimate life), feminist theory (to recover what the dominant traditions have systematically obscured), and formal methods drawn from dynamical systems theory and semiotics (to give the concept of GRW the precision it requires). The paper moves between these registers without apology, because the phenomenon demands it: a question that is simultaneously ontological, economic, psychological, political, and ethical cannot be adequately addressed from within any single disciplinary perspective.

Several chapters contain speculative formal conjectures---marked explicitly as \emph{open questions}---that go beyond what the current state of theory can rigorously establish. These conjectures are included not to claim more than the evidence supports, but because honest intellectual engagement with difficult questions requires being willing to gesture toward where the argument leads, even when the ground is not yet solid. The reader is invited to treat these sections as invitations to further inquiry rather than as settled conclusions.

\emc{A word on this paper's place in the series.} \pinkref{Paper~IX} established the spiral structure of relational generativity and the holonomy framework for distinguishing good from bad relational cycles \citep{paperix}. \pinkref{Paper~XIV} developed the eudaimonics of relational happiness, showing that happiness is the emergent signal of co-evolutionary activity in the relational field \citep{paperxiv}. The present paper asks: what is the material and symbolic infrastructure that makes this co-evolutionary activity possible? What forms of wealth support it, and what forms of wealth destroy it? The answer---that GRW is both the product and the fuel of relational co-evolution---is the political economy of the GRB framework: the account of how value is generated, stored, and circulated in the intimate relational field, and why the dominant forms of wealth in modern societies are systematically misaligned with the flourishing of intimate life.
